\documentclass[11pt,a4paper]{article}

\usepackage[ngerman]{babel}
\usepackage[utf8]{inputenc}
\usepackage[T1]{fontenc}
\usepackage{geometry}
\usepackage{microtype}
\microtypesetup{expansion=false}
\emergencystretch=3em
\usepackage{tabularx}
\usepackage{array}
\usepackage{booktabs}
\usepackage{titlesec}
\usepackage{enumitem}
\usepackage{xcolor}
\usepackage{hyperref}
\usepackage{pgfplots}
\pgfplotsset{compat=1.18}

\geometry{a4paper, margin=2.5cm}

\definecolor{accent}{RGB}{40,60,90}

\titleformat{\section}{\normalfont\Large\bfseries\color{accent}}{\thesection}{1em}{}
\titleformat{\subsection}{\normalfont\large\bfseries\color{accent}}{\thesubsection}{1em}{}

\hypersetup{
    colorlinks=true,
    linkcolor=accent,
    urlcolor=accent,
    pdftitle={Evaluation - Proof by Climbing}
}

\setlist[itemize]{leftmargin=1.5em, itemsep=0.2em}

% gemeinsame Optionen für die Balkendiagramme
\pgfplotsset{
    evalbar/.style={
        xbar,
        width=0.72\textwidth,
        height=4cm,
        xmin=0, xmax=110,
        xlabel={Zustimmung in \%},
        ytick=data,
        nodes near coords,
        nodes near coords align={horizontal},
        every node near coord/.append style={font=\small},
        bar width=10pt,
        xmajorgrids,
        major grid style={gray!25},
        axis line style={accent},
        tick label style={font=\small},
        enlarge y limits=0.3,
    }
}

\begin{document}

% ---------------- Titelseite ----------------
\begin{titlepage}
    \centering
    \vspace*{2cm}

    {\Huge\bfseries\color{accent} Proof by Climbing}\\[0.5em]
    {\Large Evaluation}\\[2cm]

    {\large Virtual Reality: Konzeption und Anwendung immersiver Technologien}\\[0.3em]

    \vfill
    {\large \today}
\end{titlepage}

% ---------------- 1. Ziel ----------------
\section{Ziel der Evaluation}

Ziel der Evaluation war es, die VR-Anwendung hinsichtlich folgender Aspekte zu bewerten:

\begin{itemize}
    \item User Experience (UX)
    \item Immersion / Präsenz
    \item VR-Komfort (Motion Sickness)
    \item Verständlichkeit der Steuerung
\end{itemize}
Die Evaluation basiert auf einem standardisierten Fragebogen, der von Testpersonen nach einem vollständigen Playthrough ausgefüllt wurde.

% ---------------- 2. Testaufbau ----------------
\section{Testaufbau}

Die Testpersonen haben einen kompletten Run im VR-System durchgeführt.
\\
Ablauf:
\begin{enumerate}
    \item Start des Spiels (Countdown: Ready -- Set -- Go)
    \item Durchführen eines vollständigen Climbing-Runs
    \item Erreichen des Ziels
    \item Ausfüllen des Fragebogens
\end{enumerate}

% ---------------- 3. Messmethodik ----------------
\section{Messmethodik}

Die Bewertung erfolgte auf einer 5-Punkt-Likert-Skala:

\begin{itemize}
    \item Stimme gar nicht zu
    \item Stimme nicht zu
    \item Neutral
    \item Stimme zu
    \item Stimme sehr zu
\end{itemize}

% ---------------- 4. Ergebnisse ----------------
\section{Ergebnisse}

\subsection{User Experience}

\begin{center}
\begin{tikzpicture}
    \begin{axis}[
        evalbar,
        symbolic y coords={Weiterempfehlung,Steuerung intuitiv,Spa{\ss} am Spiel},
    ]
    \addplot[fill=accent, draw=none] coordinates {
        (96,Spa{\ss} am Spiel) (92,Steuerung intuitiv) (100,Weiterempfehlung)
    };
    \end{axis}
\end{tikzpicture}
\end{center}
Die User Experience wurde insgesamt sehr positiv bewertet. Besonders die hohe Weiterempfehlungsrate zeigt eine starke Akzeptanz des Gameplays.

\subsection{Immersion}

\begin{center}
\begin{tikzpicture}
    \begin{axis}[
        evalbar,
        symbolic y coords={Raumgef\"uhl durch Klettern,Pr\"asenz im virtuellen Raum,Flow / vollst\"andiges Eintauchen},
    ]
    \addplot[fill=accent, draw=none] coordinates {
        (88,Flow / vollst\"andiges Eintauchen) (94,Pr\"asenz im virtuellen Raum) (90,Raumgef\"uhl durch Klettern)
    };
    \end{axis}
\end{tikzpicture}
\end{center}
Das Spiel erzeugt ein starkes Immersionsgefühl. Besonders das physikbasierte Klettern trägt zur Präsenz im virtuellen Raum bei.

\subsection{VR-Komfort (Sickness)}

\begin{center}
\begin{tikzpicture}
    \begin{axis}[
        evalbar,
        symbolic y coords={Keine Kopfschmerzen,Kein Schwindel,Keine VR-bedingten Beschwerden},
    ]
    \addplot[fill=accent, draw=none] coordinates {
        (85,Keine VR-bedingten Beschwerden) (95,Kein Schwindel) (100,Keine Kopfschmerzen)
    };
    \end{axis}
\end{tikzpicture}
\end{center}
Die VR-Anwendung ist sehr komfortabel nutzbar. Es wurden kaum Symptome von Motion Sickness festgestellt.

% ---------------- 5. Gesamtinterpretation ----------------
\section{Gesamtinterpretation}

Die Evaluation zeigt insgesamt eine sehr positive Wahrnehmung der Anwendung.
\\
Die Kombination aus physikbasiertem Klettern und minimalistischem VR-Design führt zu:

\begin{itemize}
    \item hoher Immersion
    \item guter Steuerbarkeit
    \item geringem VR-Unwohlsein
\end{itemize}

% ---------------- 6. Fazit ----------------
\section{Fazit}

Die VR-Anwendung erfüllt die Anforderungen eines immersiven und funktionalen Speedrun-Kletterprototyps.
\\
Die Testergebnisse bestätigen, dass das Gameplay sowohl verständlich als auch motivierend ist.

\end{document}
